\documentclass[conference]{IEEEtran}
\IEEEoverridecommandlockouts
% The preceding line is only needed to identify funding in the first footnote. If that is unneeded, please comment it out.
\usepackage{cite}
\usepackage{amsmath,amssymb,amsfonts}
\usepackage{algorithmic}
\usepackage{graphicx}
\usepackage{textcomp}
\usepackage{xcolor}
\usepackage[utf8]{inputenc}
\usepackage{textgreek}
\usepackage{newunicodechar}
\newunicodechar{−}{-}
\def\BibTeX{{\rm B\kern-.05em{\sc i\kern-.025em b}\kern-.08em
    T\kern-.1667em\lower.7ex\hbox{E}\kern-.125emX}}
\begin{document}
\title{AI-POWERED X-RAY VISUALIZATION AND DAMAGE DETECTION FOR PNEUMONIA\\
{\footnotesize \textsuperscript{}}
\thanks{...}
}
\author{\IEEEauthorblockN{ E.Chandralekha}
\IEEEauthorblockA{\textit{Assistant Professor,CSE} \\
\textit{Vel Tech Rangarajan}\\
\textit{Dr.Sagunthala R\&D institute}\\
\textit{of Science and Technology}\\
,Avadi, Chennai, India \\
clchandralekha08@gmail.com}
\and
\IEEEauthorblockN{ K.Nanda Kishore Reddy}
\IEEEauthorblockA{\textit{Department of CSE(AIDS)} \\
\textit{Vel Tech Rangarajan}\\
\textit{Dr.Sagunthala R\&D institute}\\
\textit{of Science and Technology}\\
Avadi, Chennai, India \\
vtu21896@veltech.edu.in}
\and
\IEEEauthorblockN{N.Venu Gopal Reddy  }
\IEEEauthorblockA{\textit{Department of CSE(AIDS)} \\
\textit{Vel Tech Rangarajan}\\
\textit{Dr.Sagunthala R\&D institute}\\
\textit{of Science and Technology}\\
Avadi, Chennai, India \\
vtu21798@veltech.edu.in}
}
\maketitle
\begin{abstract}
Pneumonia still kills millions each year, especially those who are most vulnerable. Deep learning has made it easier to tell pneumonia from chest X-rays, but many of the latest systems are still leaving doctors scratching their heads. They're accurate enough, but don't give a window into the lungs to give a clue about what's happening, nor help explain its severity. And that's where our research comes in. We've built an AI that combines a DenseNet121 classifier colourises the resultant 3D lungs, and gives a grading for severity! It's not just a quick in, throw your X-ray at the system, and out pops a diagnosis (achieving 85.26\% accuracy on the Kaggle Chest X-Ray Pneumonia dataset). It highlights regions of concern, produces colour-coded 3D images for quick reference, and distributes automatically graded staging of disease severity. A web-based interface is also provided for clinical deployment. In testing, the model achieved 85.26\% accuracy, 87.25\% weighted precision, 85.26\% weighted recall, 84.39\% weighted F1-score, and 0.96 AUC-ROC. Beyond quantitative metrics, the system's visual explanations enhance transparency and support clinician trust in AI-assisted diagnosis.
\end{abstract}
\begin{IEEEkeywords}
Deep Learning, Explainable AI, Grad-CAM, 3D Visualization, Pneumonia Detection, Chest X-Ray, DenseNet121, Medical Imaging
\end{IEEEkeywords}


\section{\textbf{INTRODUCTION}}
\noindent\hspace{0.5cm}Pneumonia is a disease of the lungs that is responsible for millions of respiratory sicknesses across the world. In fact, it's so serious that recent studies show it is responsible for about 15\% of childhood deaths worldwide, killing some 2.5 million children under 5 every year. Detecting the disease early makes all the difference in the effectiveness of treatment. But here's the kicker: it is really hard to read X-rays. It's a skill, and doctors don't even agree often, with disagreements in reading that are sometimes as high as 70-85\%. Doctors primarily if not exclusively use chest X-rays for the diagnosis of pneumonia. It's quick, cheap, and does not expose the patient to heavy doses of radiation. Unfortunately, X-rays are flat images and many of the signs of Pneumonia are subtle, so it is sometimes difficult to judge with certainty what you're looking at. For all these reasons, there is a move on to have computer-aided diagnosis (CAD) systems to aid and abet clinical decisions. This work addresses these problems. It proposes a robust framework for automatic pneumonia detection using X-ray images based on DenseNet121 classification tuned with transfer learning, and uses Grad-CAM for activation mapping as well as state-of-theart projection of activation on a parametric surface to create a 3D model of the lung region. It can automatically stage pneumonia severity as mild, moderate or severe,and also provide a report in a ready-to-deploy web application.

\section{\textbf{LITERATURE SURVEY}}
\hspace{0.5cm}Detecting and visualising pneumonia with AI is on the rise as there's a pressing demand for accurate medical imaging. Researchers are Resorting to deep learning models like DenseNet, ResNet and VGG16 to reliably and accurately diagnose pneumonia from chest x-rays. DenseNet, which has a ton of connections, can pass features to each other and capture detailed information. ResNet, however, derives its depth from a technique called residual learning, meaning it can discover those messy, complex patterns in the lungs. There's also VGG16 which is deep in terms of its layers, allowing it to exploit hierarchical configuration for both spatial as well as textural features – crucial for localizing where the infection is. And of course it's not just about getting the answer right. Visualization tool, Grad-CAM comes into play to show which part of the lung the model focuses on, providing clinicians insight and faith into these AI systems. A huge stride in the name of making AI accurate yet transparent and applicable in real-world settings.\\
Deep learning for pneumonia detection has gained significant traction recently, and hospitals sure need quicker and more accurate ways of interpreting medical images. Rajpurkar and his colleagues developed CheXNet, which is a deep convolutional neural network based on the DenseNet architecture, and found that CheXNet performed on par with radiologists in identifying pneumonia on chest X-rays. "We show that deep learning for pneumonia detection could assist doctors in providing a diagnosis that is accurate and reliable[1].\\
Kermany and colleagues made a deep learning model analyzing medical images for pneumonia. They show AIs that actually can support a doctor's recognition of patterns in images that are somewhat subtle, showing good improvement in early detection across several datasets. This means not only in the lab, they also validated the technique in a medical environment. [2].\\
Liang and Zheng do pediatric pneumonia detection using transfer learning on a deep residual network. They show that a pre-trained deep model can learn to do medical imaging, even if they do not have that much pediatric data. That is the performance high and diagnosis is more reliable using relatively few images. [3].\\
Stephen et al. developed a deep learning framework that accelerates pneumonia classification. They concentrated on optimizing CNNs to operate faster and achieve higher precision, allowing doctors to receive real-time decision support when it matters most [4].\\
Sharma and colleagues use a hybrid deep learning approach using VGG-16 and added neural layers to strengthen pneumonia detection in chest X-ray images. Their tuned CNNs picked up features strongly representative of accidently detected pneumonia in spatial information, producing strong CNNs capable of being accurately tested against others. [5].\\
Selvaraju et al. proposed a gradient-based method called Grad-CAM that tells us what parts of the image a CNN uses to make its decision; useful for pneumonia detection as it shows where the infected lung is, and understanding AI's reasoning behind its predictions [6].\\
Guan and Huang developed multi-label chest X-ray classification with category-wise residual attention learning; simply, knowing where to look—concentrating on those pixels in the X-ray that matter for detecting disease—gives greater accuracy as well as adequate explanations on how the AI see pneumonia [7].\\
Wang et al. created ChestX-ray8, the largest publicly available dataset of more than 100,000 labelled X-ray images covering eight different thoracic diseases; routinely used as a benchmark to train and evaluate AI models with weak supervision for disease classification and localization, including pneumonia [8].\\
Chen et al. contributed by applying deep learning models to 3D reconstruction of lungs from 2D chest x-rays; it resonates with our approach providing a new perspective on medical reconstruction with glimpses into three dimensions of the structure and damage to lung, an insight into the spread and severity of pneumonia [9].\\
Kermany and Goldbaum created a labelled public dataset of chest X-rays and OCT, quickly making it a favourite for AI researchers to train and test AI models for pneumonia detection; being reliable and open-source has helped standardisation and acceleration in medical AI research [10]..\\
Huang and et al made DenseNet, "a neural net with access from every layer to every other layer in a feed-forward fashion, rather than just to the previous layer in the series.This lets the model reuse features more easily, and the learning becomes a lot easier and more stable due to the flow of information between sparely connected layers rather than "chained in" one after the other, so the convolutional networks can be stronger and much tighter. Some medical imaging tools use DenseNet as their backbone, like pneumonia detection tools like CheXNet.[11]\\
Irvin et al. create a huge dataset of chest X-ray images as well as providing 14 observations per each labeled example along with uncertainty values. They demonstrate that when a model learns what label uncertainty is, earlier diagnostic performance improves, and they counterpose their models with human expert radiologists. This dataset is used to form a robust classifier of diseases detectable in chest X-rays.[12].\\
Lakhani and Sundaram use convolutional neural networks for the detection of pulmonary tuberculosis from chest radiographs. They compare their CNN to traditional feature based algorithms and, among other results, conclude that some deep CNNs are capable of highest diagnostic accuracy. They set the scene for deep learning of chest X-ray−based pulmonary disease classification.[13].\\
Zhou et al. introduce Class Activation Mapping (CAM), a method that allows weakly supervised localization by displaying those parts of the image discriminative for the class. By doing this, they improve model interpretability by showing that part of the convolutional features lead logically to class prediction, with no need for an annotated bounding-box approach. This concept forms the basis of other medical imaging explainability methods later on.[14].\\
Ye et al. study pneumonia detection in chest X-ray images using deep learning with pretrained CNNs. Again showing how deep learning can speed up conventionally slow training times, and improve convergence, showing it adequate for medical image classification.[15].\\

\section{\textbf{METHODOLOGY}}
\hspace{0.5cm} Here's how the Pneumonia Detection System works, step by step. The team set out to classify chest X-rays either normal or showing signs of pneumoni and mapped out the process in a system architecture diagram.

This diagram contains all the steps, starting with data preparation, model training, and model deployment up to the interpretation of the output of the model. Firstly, we gathered some images of chest X-rays of a Kaggle collection of real medical cases, all of which were either annotated as healthy or infected with pneumonia. In this manner of sampling, we were able to have a balanced dataset, that is, one was made up of equal numbers of normal and pneumonia. They applied these labels with the purpose of training the model, as well as testing the model performance Next stop, pre-processing. All X-ray images were downsized to an equal number of 224×224 pixels, thus, all the data fitted together in the model. We also scaled the pixel values, to make the model learn more aid in and curb problems by variations in acquisition of the images. We used small scale data augmentation methods to train the model to distinguish between good and bad pictures such as random flipping and rotating of images. This step enhanced the variety of the dataset and made the model more resistant to the process of identifying pneumonia in X-rays that were not seen.
\vspace{0.5 cm}
\begin{figure}[h]
    \centering
    \includegraphics[width=1.0\linewidth]{xray architecture.png}
    \caption{Architecture Overview}
    \label{fig:enter-label}
\end{figure}
\\
First off, there is the chest X-ray itself that undergoes a resize, normalization, alignment/registering and is processed by the DenseNet121 model, which does the heavy lifting to poke around the image, determining the presence or absence of pneumonia. Upon making a call, the system will have two explainability pathways. One of them is Grad-CAM that generates a heatmap image of the X-ray to identify suspicious spots in 2-D. Meanwhile, a 3D lung model indicates the locations of these activations on a parametric lung surface indicating where exactly the model believes the trouble spots to be.
\\
At this point, a 2D heatmap and 3D lung visualization can be displayed to identify the level of severity of the situation: normal, mild, moderate, and severe. It is important to note that the severity staging is derived from the Grad-CAM activation coverage and intensity rather than radiologist-annotated ground-truth severity labels; this heuristic provides a clinically suggestive but not externally validated indicator of disease extent. In addition to creating a simple prediction the system will automatically generate a clinical report that the model describes its results, visually explains the most significant areas, and gives a general severity rating.
The pneumonia detection system begins with first preprocessing a chest x-ray by rescaling, normalizing the x-ray, and converting the image. After that, the processed picture is transferred to DenseNet121, which determines whether the scan is normal or is pneumonia-affected. When classified, two visual explanations are produced. Grad-CAM generates a 2D heatmap of areas of the lungs that contributed to the prediction, and a 3D visualization plots the intensity of the activation on a lung model to give the spatial location of pneumonia. According to these visual outputs, the system then estimates disease severity and assigns a level, that is, normal, mild, moderate and severe, based on the proportion and intensity of Grad-CAM activations. Since the underlying Kaggle dataset provides only binary labels (Normal vs.\ Pneumonia) and no radiologist-graded severity annotations, the severity levels reported by the system represent model-derived estimates and have not been validated against clinical ground truth. It will then automatically produce a clinical report with the prediction, level of severity, heatmap, 3D lung visualization, and key observations and thus the output can be understood easily and can be used clinically. Fig. 1 shows the general architecture of the proposed pneumonia detection system, which shows the preprocessing pipeline, DenseNet121 backbone, explainability modules, and reporting workflow.
\\
\subsection{\textbf{Dataset Description}}l
The data employed in the project is a chest X-ray data set which is subdivided into two groups; Pneumonia and Normal. The data were acquired in the Kaggle Chest X-ray repository and backed with the records of trusted clinical resources to provide reliability to the data. The normal group consists of normal chest X-rays and the Pneumonia group consists of both pediatric and adult cases of clinically verified bacterial and viral. The dataset has a total of 5,863 images, 1,583 of them are Normal and 4,280 are Pneumonia images.
Pretrained convolutional neural networks were used; therefore, all images were resized to 224 × 224 pixels and three-channel RGB format was obtained. The data were divided into training, validation and testing sets following the Kaggle repository's default partitioning. The validation set comprises only 16 images (8 per class), while the test set contains 624 images (234 Normal and 390 Pneumonia). Due to the extremely small validation split, validation metrics should be interpreted with caution, and the held-out test set provides the definitive performance evaluation.

%------------------------------------------------------------

\subsection{\textbf{Preprocessing}}
Before training the model, I ran a few preprocessing steps to make learning smoother and help the model generalize better.
\subsubsection{Normalization}
First off, I normalized the pixel values so everything fell between 0 and 1:
\begin{equation}
I_{norm} = \frac{I - I_{min}}{I_{max} - I_{min}}
\end{equation}
where \( I \) represents the original pixel intensity.
\subsubsection{Data Augmentation}
To keep the model from overfitting and to make it more robust, I mixed in some data augmentation tricks during training:
\begin{itemize}
    \item Randomly flipped images left to right
    \item Rotated images by up to 10 degrees in either direction
    \item Tweaked the contrast a bit
    \item Added a little zoom
\end{itemize}
\subsubsection{Batch Generation}
I set up mini-batches of 32 images each. This kept the GPU running efficiently and made sure gradient updates stayed stable while training.

%------------------------------------------------------------

\subsection{Model Architecture}
\subsubsection{DenseNet121 Classifier}
Fig. 2 demonstrates the tailored DN121 architecture to pneumonia classification with an accent on the adjusted layers and the fine-tuning approach.
DenseNet121 was selected as the backbone architecture because its dense connections enable effective feature reuse and support smooth gradient propagation. The original classification head was replaced, and a new custom classification head was designed.
\begin{itemize}
    \item Global Average Pooling layer
    \item Fully connected layer with 512 neurons and ReLU activation
    \item Dropout layer with a 0.3 rate
    \item Final dense layer with Softmax activation for binary classification
\end{itemize}
\begin{figure}[h]
    \centering
    \includegraphics[width=0.95\linewidth]{dense net 121.png}
    \caption{Customized DenseNet121 architecture for pneumonia classification}
    \label{fig:densenet}
\end{figure}
\begin{table}[h]
    \centering
    \caption{Model Evaluation Metrics for DenseNet121}
    \label{tab:model_comparison_pneumonia}
    \begin{tabular}{|l|c|c|c|c|c|}
    \hline
    \textbf{Class} & \textbf{Precision} & \textbf{Recall} & \textbf{F1-Score} & \textbf{Support} \\
    \hline
    \textbf{Normal} & 0.97 & 0.63 & 0.76 & 234 \\
    \hline
    \textbf{Pneumonia} & 0.82 & 0.99 & 0.89 & 390 \\
    \hline
    \textbf{Macro Avg} & 0.89 & 0.81 & 0.83 & 624 \\
    \hline
    \textbf{Weighted Avg} & 0.87 & 0.85 & 0.84 & 624 \\
    \hline
    \end{tabular}
\end{table}
The results of the DenseNet121 model, in terms of accuracy, precision, recall, F1-score, and AUC–ROC, utilized to determine the classification result, are displayed in Table I.

%------------------------------------------------------------

\subsubsection{Grad-CAM Based Explainability}

We used Grad-CAM to make the predictions of our model interpretable. In simple terms, Grad-CAM tells us which parts of the lungs had the greatest influence on our model's predictions. For each feature map, we derive an importance weight, and then give us a heatmap that conveys which part of the lung the model was focusing on, and, in our case, the infected part of the lung. This makes our results more interpretable, believing what our AI is seeing is:
\begin{equation}
\alpha_k^c = \frac{1}{Z} \sum_{i}\sum_{j} \frac{\partial y^c}{\partial A_{ij}^k}
\end{equation}

The Grad-CAM heatmap is generated as:
\begin{equation}
L_{GradCAM}^c = \text{ReLU}\left(\sum_k \alpha_k^c A^k\right)
\end{equation}

%------------------------------------------------------------

\subsubsection{3D Lung Visualization}

To get a better sense of where the infection is inside the lungs, we took the Grad-CAM heatmaps and projected them onto a 3D model of the lungs using Plotly. By mapping activation intensity to different colors, you can quickly spot which areas are more severely affected. It should be noted that this 3D representation is generated by projecting the 2D Grad-CAM heatmap onto a generic parametric lung surface rather than reconstructing actual patient-specific 3D anatomy from the X-ray. It therefore serves as an approximate spatial visualization to aid clinical interpretation, not as a true volumetric reconstruction.

%------------------------------------------------------------

\subsection{Training Configuration}
For training, we used PyTorch. Adam was our optimizer because it adapts well during learning. The parameter updates followed Adam's standard rule, with η as the learning rate for the bias-corrected moment estimates. We relied on categorical cross-entropy for the loss function.
\begin{equation}
\theta_{t+1} = \theta_t - \eta \frac{\hat{m}_t}{\sqrt{\hat{v}_t} + \epsilon}
\end{equation}
where \( \eta \) denotes the learning rate, and \( \hat{m}_t \) and \( \hat{v}_t \) are bias-corrected first and second moment estimates.

The categorical cross-entropy loss function was employed:
\begin{equation}
\mathcal{L} = -\sum_{i=1}^{C} y_i \log(\hat{y}_i)
\end{equation}

To keep the model from overfitting and to help it learn smoothly, we used early stopping and ReduceLROnPlateau scheduling..

%------------------------------------------------------------

\section{Result Analysis}

We designed and implemented an AI-driven Pneumonia Detection System using chest X-rays; some have healthy lungs and some are with pneumonia. Our modeld for the backbone is DenseNet121, used in conjunction with Grad-CAM and 3D lung visualization for explainability of the predictions. We split our data 80\% for training, and 10\% to both validate and testing.

\subsection{Performance Evaluation}

In the case of performance indicators, we came up with the common medical AI measures; accuracy, F1-score, precision and recall. That provides the model with just how it actually selected pneumonia not in a general sense, but in interesting ways clinically.

\subsubsection{Accuracy}
\begin{equation}
\text{Accuracy} = \frac{TP + TN}{TP + TN + FP + FN}
\end{equation}

\subsubsection{Precision}
\begin{equation}
\text{Precision} = \frac{TP}{TP + FP}
\end{equation}

\subsubsection{Recall}
\begin{equation}
\text{Recall} = \frac{TP}{TP + FN}
\end{equation}

\subsubsection{F1-Score}
\begin{equation}
\text{F1-Score} = \frac{2 \times \text{Precision} \times \text{Recall}}{\text{Precision} + \text{Recall}}
\end{equation}
Such metrics are the solution of the entire story as they are both false positive and negative - X-ray diagnostics cannot risk giving such incorrect diagnosis. The fine-tuned DenseNet121 achieved a classification accuracy of 85.26\% on the held-out test set. The dense connectivity pattern of DenseNet121 enables effective feature reuse, contributing to robust classification performance. Grad-CAM produced heatmaps of the pneumonia-affected lung regions that provided clinicians with confidence in the model's observations. Constant precision and management of novel incidences to other age categories of patients. The training and validation accuracy achieved in the course of learning is summarized in Table II and it shows convergence of the model. Note that the validation set comprises only 16 images from the Kaggle dataset's default split, and therefore validation accuracy values are highly variable and should not be treated as a reliable generalization metric; the test set results in Table I provide the definitive performance evaluation.\\
\begin{table}[h]
\centering
\caption{Training and Validation Accuracy}
\label{tab:epoch_accuracy}
\begin{tabular}{|c|c|c|}
\hline
\textbf{Epoch} & \textbf{Train Accuracy (\%)} & \textbf{Validation Accuracy (\%)} \\ \hline
1  & 95.11 & 100.00 \\ \hline
2  & 97.18 & 87.50 \\ \hline
3  & 97.55 & 93.75 \\ \hline
4  & 97.97 & 75.00 \\ \hline
5  & 98.10 & 100.00 \\ \hline
6  & 98.59 & 100.00 \\ \hline
7  & 98.70 & 87.50 \\ \hline
8  & 98.95 & 100.00 \\ \hline
9  & 98.87 & 100.00 \\ \hline
10 & 98.85 & 100.00 \\ \hline
11 & 99.08 & 100.00 \\ \hline
12 & 99.29 & 93.75 \\ \hline
13 & 99.52 & 81.25 \\ \hline
14 & 99.19 & 100.00 \\ \hline
15 & 99.25 & 100.00 \\ \hline
16 & 99.42 & 100.00 \\ \hline
17 & 99.58 & 100.00 \\ \hline
18 & 99.58 & 100.00 \\ \hline
19 & 99.48 & 100.00 \\ \hline
20 & 99.54 & 100.00 \\ \hline
\end{tabular}
\end{table}

\noindent\textbf{2. Training Performance Analysis:}\\
It was a fast learner, and was stable for all 20 epochs. It hit a training score of 95.11\% from the get-go and ended up at 99.54\%, which indicates that it focused on the right features and continued to improve. Validation accuracy ranged from 75\% to 100\% across epochs; however, the validation set consists of only 16 images (8 per class) as provided by the Kaggle dataset's default partition. With such a small validation set, each misclassified image shifts accuracy by 6.25 percentage points, making these fluctuations statistically unreliable as a measure of generalization. The held-out test set of 624 images provides a more dependable evaluation of model performance.
Training loss was rapidly decreased from 0.1373 to 0.0108. Validation loss was also low. This indicates that the model successfully reduced errors during training. One reassuring aspect is the consistent improvement of training accuracy throughout training, indicating that the model progressively learned relevant features. Solid augmentation strategies, dropout, learning rate scheduling, and early stopping came together for good effect.
From a clinical aspect, these results imply that it can be considered a useful tool to distinguish pneumonia from normal Chest X Rays. The model tends to err on the side of detecting pneumonia, as reflected in its high pneumonia recall of 98.72\%, which is a preferable bias in a screening context where missing a positive case carries greater risk than a false alarm. In summary, the training results should convey the message that this AI system is stable and functional, though further evaluation on larger and more diverse validation sets would strengthen the generalization claims. Fig. 3 demonstrates the general performance measures of the model, which gives a comparative representation of the accuracy, precision, recall, and F1-score.\\
\begin{figure}[h]
    \centering
    \includegraphics[width=0.9\linewidth]{model performance evaluation.png}
    \caption{Model Performance Metrics}
    \label{fig:model_perf}
\end{figure}
\\
\noindent\textbf{3. Model Performance Metrics:}\\
Performance of the pneumonia classifier was tested and recorded as 85.26\%, meaning that a greater percentage of the test images have been classified appropriately. The precision for the pneumonia class was 81.57\%, indicating that the majority of the cases predicted as pneumonia were true positives. The value of recall was 98.72\%, which was very high, indicating that most actual cases of pneumonia were correctly identified with very few cases missed. The F1-score for the pneumonia class was 89.33\%, depicting a reasonable balance between false alarms and sensitivity. The weighted precision across both classes was 87.25\%, and the weighted F1-score was 84.39\%. It is worth noting that the model's high recall comes at the cost of lower precision for the normal class (specificity of 62.82\%), meaning a notable fraction of normal cases are misclassified as pneumonia.

The AUC-ROC measure was 95.77\% which indicates a strong ability to differentiate between normal and pneumonia cases in the chest X-rays. Judging by these findings, the model may be deemed as suitable for pneumonia screening with the caveat that the relatively lower specificity should be addressed in future iterations. A limitation of the current evaluation is the absence of direct baseline comparisons with alternative architectures (e.g., ResNet50, VGG-16, or EfficientNet) trained on the same dataset; such comparisons would help contextualize DenseNet121's relative performance. Fig. 4 shows the analysis of ROC curve, which shows that the model was able to differentiate between normal and pneumonia-infected X-ray images of the chest. Fig. 5 also shows the ROC curve analysis at different thresholds that would attest to the same discriminative performance being generated at different operating points.\\
\vspace{0.5 cm}
\begin{figure}[h]
    \centering
    \includegraphics[width=0.9\linewidth]{training_history (2).png}
    \caption{ROC Curve Analysis}
    \label{fig:roc_curve1}
\end{figure}
\begin{figure}[h]
    \centering
    \includegraphics[width=0.9\linewidth]{training_history (1).png}
    \caption{ROC Curve Analysis}
    \label{fig:roc_curve2}
\end{figure}
\\
\noindent\textbf{4. ROC Curve Analysis:}\\
Evaluation of the capacity of the model to differentiate actual cases was assessed by the use of an ROC curve. The curve is a lot higher than the diagonal, which means that the performance is much higher than random guessing. The obtained AUC was 95.77\% which represents very good discriminative power on a balanced dataset. The sensitivity, compared to the false-positive rate, is high proving that, pneumonia cases, and normal cases are successfully discriminated.

Such an outcome has clinical implications because various medical settings demand various trade-offs between sensitivity and specificity such as mass screening and emergency triage. The ROC curve also reveals that the constant performance is ensured at different decision threshold levels, and high true positive rate is preserved when other parameters are varied. As a result, the model is reliable in a variety of clinical settings, such as mass screening and situations with a time constraint of diagnosis.
\vspace{0.5 cm}
\begin{figure}[h]
    \centering
    \includegraphics[width=0.5\linewidth]{Grad-CAM Heatmap Analysis.jpg}
    \caption{Grad-CAM Heatmap Analysis}
    \label{fig:gradcam}
\end{figure}
\\
\begin{figure}[h]
    \centering
    \includegraphics[width=0.7\linewidth]{lung output normal.png}
    \caption{3D Visualization Normal Lungs}
    \label{fig:3d_normal}
\end{figure}
\begin{figure}[h]
    \centering
    \includegraphics[width=0.6\linewidth]{lung output pnemonia.png}
    \caption{3D Visualization Pneumonia Lungs}
    \label{fig:3d_pneumonia}
\end{figure}

\noindent\textbf{5. Visualization and Output:}\\
To make the model's reasoning transparent, we take the input X-ray and generate a Grad-CAM heatmap to show where the model focuses most strongly when making its prediction. The hotter areas indicate where the model believes infection is present. This provides clinicians with interpretable evidence supporting the predictions.
Beyond Grad-CAM, the system also generates a 3D surface of the lungs by projecting the activation intensities onto a parametric lung model. Clinicians can see not just how severe the pneumonia might be, but where it is spatially located. The 3D rendering provides additional spatial context for location and estimated severity. As noted earlier, since the 3D model is a projection of 2D activations onto a generic surface rather than a patient-specific reconstruction, it should be interpreted as an approximate visualization aid. The final output of the model is the predicted class (Normal or Pneumonia), the confidence score, the Grad-CAM heatmap, and also the 3D presentation of the lung with the activation strength projected onto the anatomy to indicate the regions of concern. Fig. 6 illustrates the Grad-CAM heatmap analysis, which identifies the lung areas that had the greatest impact on the model's pneumonia prediction. Fig. 7 demonstrates the 3D image of normal lungs, the image indicates the patterns of baseline activation of the healthy cases of the chest X-ray. Fig. 8 shows the 3D visualization of lungs with pneumonia which mapped the high activation intensities to spatial regions related to the pathological outcomes.\\



\section{\textbf{CONCLUSION AND FUTURE WORK}}
This article proposes an AI-based pneumonia detection system that identifies the chest X-rays as Normal or Pneumonia by the use of a fine-tuned DenseNet121 backbone, with the assistance of consistent preprocessing, effective training, and in-depth validation to reach a high level of predictive and clinical accuracy. Grad-CAM heatmaps that display the lung locations of diagnostic interest in a clear way and a 3D lung surface projection that provides a sense of space, essential to surmounting black-box constraints, make the system suitable in applications to screening and teleradiology. The system is ready to be used in clinical practice as it combines proper classification, visual explainability, and severity-oriented reporting. Several limitations should be acknowledged: the severity staging (mild, moderate, severe) is derived from Grad-CAM activation heuristics and has not been validated against radiologist-annotated severity labels; the 3D visualization projects 2D heatmaps onto a generic lung surface rather than reconstructing actual patient anatomy; and no direct baseline comparisons with alternative architectures were performed on the same dataset. Future work can address these limitations and expand this framework to multi-disease thoracic screening with baseline comparisons against ResNet, VGG-16, and EfficientNet, add patient metadata to provide contextual risk assessments, validate the severity staging against radiologist ground truth on appropriately annotated datasets, enable smooth DICOM/PACS and clinical workflow integration, facilitate edge and real-time inference, provide explainability with uncertainty and concept-level attributions, add longitudinal severity tracking, and multimodal imaging (CT and ultrasound) and interactive report generation enabling clinicians to refine, finalize, and export results directly to electronic medical record systems.



\bibliographystyle{IEEEtran}
\begin{thebibliography}{99}

\bibitem{rajpurkar2017}
P.~Rajpurkar \textit{et al.}, ``CheXNet: Radiologist-Level Pneumonia Detection on Chest X-Rays with Deep Learning,'' \emph{arXiv preprint arXiv:1711.05225}, 2017.

\bibitem{kermany2018}
D.~S.~Kermany \textit{et al.}, ``Identifying Medical Diagnoses and Treatable Diseases by Image-Based Deep Learning,'' \emph{Cell}, vol. 172, no. 5, pp. 1122–1131, Feb. 2018.

\bibitem{liang2020}
G.~Liang and L.~Zheng, ``A Transfer Learning Method with Deep Residual Network for Pediatric Pneumonia Diagnosis,'' \emph{Computer Methods and Programs in Biomedicine}, vol. 187, p. 104964, Apr. 2020.

\bibitem{stephen2019}
O.~Stephen, M.~Sain, U.~J.~Maduh, and D.~U.~Jeong, ``An Efficient Deep Learning Approach to Pneumonia Classification in Healthcare,'' \emph{Journal of Healthcare Engineering}, vol. 2019, Article ID 4180949, 2019.

\bibitem{sharma2021}
S.~Sharma, K.~Guleria, S.~Tiwari, and S.~Kumar, ``A Deep Learning Based Model for the Detection of Pneumonia from Chest X-Ray Images Using VGG-16 and Neural Networks,'' \emph{Procedia Computer Science}, vol. 218, pp. 357–366, 2021.

\bibitem{selvaraju2017}
R.~R.~Selvaraju, M.~Cogswell, A.~Das, R.~Vedantam, D.~Parikh, and D.~Batra, ``Grad-CAM: Visual Explanations from Deep Networks via Gradient-Based Localization,'' in \emph{Proc. IEEE International Conference on Computer Vision (ICCV)}, Venice, Italy, Oct. 2017, pp. 618–626.

\bibitem{guan2018}
Q.~Guan and Y.~Huang, ``Multi-Label Chest X-Ray Image Classification via Category-Wise Residual Attention Learning,'' \emph{Pattern Recognition Letters}, vol. 130, pp. 259–266, Feb. 2020.

\bibitem{wang2020}
X.~Wang \textit{et al.}, ``ChestX-Ray8: Hospital-Scale Chest X-Ray Database and Benchmarks on Weakly-Supervised Classification and Localization of Common Thorax Diseases,'' in \emph{Proc. IEEE Conference on Computer Vision and Pattern Recognition (CVPR)}, Honolulu, HI, USA, Jul. 2017, pp. 2097–2106.

\bibitem{chen2022}
H.~Chen, Y.~Zhao, L.~Zhang, and W.~Wang, ``3D Lung Reconstruction from Chest X-Ray Using Deep Learning,'' \emph{Medical Image Analysis}, vol. 78, p. 102401, May 2022.

\bibitem{kermany2018dataset}
D.~S.~Kermany and M.~Goldbaum, ``Labeled Optical Coherence Tomography (OCT) and Chest X-Ray Images for Classification,'' \emph{Mendeley Data}, v2, 2018. [Online]. Available: http://dx.doi.org/10.17632/rscbjbr9sj.2

\bibitem{huang2017}
G.~Huang, Z.~Liu, L.~Van der Maaten, and K.~Q.~Weinberger, ``Densely Connected Convolutional Networks,'' in \emph{Proc. IEEE Conference on Computer Vision and Pattern Recognition (CVPR)}, Honolulu, HI, USA, Jul. 2017, pp. 4700–4708.

\bibitem{irvin2019}
J.~Irvin \textit{et al.}, ``CheXpert: A Large Chest Radiograph Dataset with Uncertainty Labels and Expert Comparison,'' in \emph{Proc. AAAI Conference on Artificial Intelligence}, Honolulu, HI, USA, Jan. 2019, pp. 590–597.

\bibitem{lakhani2017}
P.~Lakhani and B.~Sundaram, ``Deep Learning at Chest Radiography: Automated Classification of Pulmonary Tuberculosis by Using Convolutional Neural Networks,'' \emph{Radiology}, vol. 284, no. 2, pp. 574–582, Aug. 2017.

\bibitem{zhou2016}
B.~Zhou, A.~Khosla, A.~Lapedriza, A.~Oliva, and A.~Torralba, ``Learning Deep Features for Discriminative Localization,'' in \emph{Proc. IEEE Conference on Computer Vision and Pattern Recognition (CVPR)}, Las Vegas, NV, USA, Jun. 2016, pp. 2921–2929.

\bibitem{ye2020}
J.~Ye, T.~Cheng, Z.~Yi, and Z.~Zhang, ``Pneumonia Detection and Classification Based on Transfer Learning in X-Ray Images,'' \emph{Frontiers in Public Health}, vol. 8, p. 146, Apr. 2020.

\end{thebibliography}

\end{document}
